\section{Anexo}

\subsection{Conversión analógico-digital}

La tensión $V_{out}$ generada por el divisor resistivo es una señal analógica. Para que pueda ser procesada por el Arduino, esta tensión es convertida a un valor digital mediante el conversor analógico-digital (ADC).

El Arduino UNO utiliza un ADC de 10 bits en sus entradas analógicas. Por lo tanto, la conversión dispone de:

\begin{equation}
2^{10}=1024
\end{equation}

niveles de cuantización, correspondientes a valores digitales entre 0 y 1023.

Considerando una tensión de referencia $V_{ref}$, la tensión aplicada a la entrada analógica puede expresarse idealmente como:

\begin{equation}
V_{out}
=
\frac{ADC}{1023}V_{ref},
\end{equation}

donde $ADC$ representa el valor digital obtenido por el conversor.

En el circuito utilizado, la tensión de alimentación es de $5~\mathrm{V}$ y se utiliza como referencia para la medición, por lo que:

\begin{equation}
V_{out}
=
\frac{ADC}{1023}5~\mathrm{V}.
\end{equation}

Combinando esta expresión con la ecuación del divisor resistivo, se obtiene:

\begin{equation}
R_{NTC}
=
R
\left(
\frac{1023}{ADC}-1
\right).
\end{equation}

Esta expresión permite calcular directamente la resistencia del termistor a partir del valor obtenido por la entrada analógica del Arduino. La ecuación coincide con la utilizada en el programa proporcionado por la cátedra para obtener la resistencia del NTC a partir de la lectura del ADC. 
